% ================================================================
% SECTION VII: DISCUSSION
% ================================================================

\subsection{Trade-offs Between the Three Paradigms}
\label{sec:disc-paradigms}

\paragraph{DAS protocols (Celestia, Avail, Ethereum)}
These give light clients the strongest guarantees, because a resource-constrained node can check availability itself.
The guarantee weakens if too few nodes sample, a dependence the other two paradigms do not have.

\paragraph{Validator-recovery protocols (Polkadot, NEAR)}
Polkadot's finality gate (GRANDPA will not finalize blocks without sufficient availability bitfields) is the most direct enforcement we observe: unavailable data cannot be finalized.
In exchange, light clients cannot verify availability themselves.
For parachain validation this is acceptable, because the relay chain validators are the main consumers of the data; an external rollup using Polkadot or NEAR for DA would have to trust the validator set.

\paragraph{VID protocol (Espresso)}
None of the other protocols has an equivalent of Tiramisu's algebraic bribery resistance, but in Mainnet~0 it counts for little while the operator set is permissioned (20~operators) and there is no slashing.
The 1~wei/byte fee is likewise a placeholder.
Whether VID's advantages matter in practice will depend on a permissionless deployment with economic security behind it.

\subsection{The Cost--Throughput Frontier}
\label{sec:disc-frontier}

Fig.~\ref{fig:landscape}(a) plots median cost against protocol maximum throughput, and there is no simple inverse relationship.
Espresso combines low cost with moderate throughput, but its cost is a Mainnet~0 placeholder.
Celestia and Avail have moderate throughput at moderate cost.
Polkadot has the highest theoretical throughput ($\sim$167~MiB/s) at $\sim$\$5$\times$10$^{-6}$/MiB, although coretime pays for execution as well as DA.
NEAR combines high capacity ($\sim$59~MiB/s) with the highest cost among active fee markets ($\sim$\$0.53/MiB), the opposite corner from Polkadot.
After Fusaka, Ethereum offers modest capacity at a cost comparable to Celestia and Avail, with the largest validator set and stake behind it.

The protocols were built for different jobs: Polkadot for parachain validation, NEAR for general smart-contract execution, Ethereum for censorship resistance and economic security, and Celestia and Avail as dedicated DA layers.
Cost and throughput therefore do not rank them on their own; retention, trust assumptions, light-client verification and ecosystem integration matter as much.

\subsection{Polkadot's Scaling Headroom}
\label{sec:disc-polkadot}
With only $\sim$26 of 100 cores engaged (Fig.~\ref{fig:blocks}(c)), Polkadot has $\sim$74\% idle DA capacity: throughput could grow from the current upper bound of $\sim$43~MiB/s toward $\sim$167~MiB/s without parameter changes.
As engagement rises, the coretime market should raise core prices; at the current $\sim$\$5$\times$10$^{-6}$/MiB, cost is unlikely to limit near-term growth.
Because actual PoV sizes are not exposed on-chain, a more precise capacity analysis requires PoV-level data.

\subsection{Implications for Rollup Developers}
\label{sec:disc-rollups}

Retention constrains a rollup's choice of DA layer before cost does.
An optimistic rollup with a 7-day challenge period cannot rely on Polkadot (25~h) or NEAR ($\sim$1.5~d) without its own archive, and Celestia's 7-day sampling window leaves no margin.
For the same reason, spot prices mislead on short-retention layers; the annualized cost (Section~\ref{sec:meth-metrics}) counts the reposts needed to keep data available.
Fee floors and token prices, more than congestion, currently set what rollups pay: Ethereum's post-Fusaka cost sits near the EIP-7918 reserve price, and Celestia's largest demand surge did not raise its median fee.
With every protocol below 40\% mean utilization, capacity does not yet separate the options; how fee markets behave under load will matter once utilization rises.

\subsection{Protocol Design Observations}
\label{sec:disc-design}

\paragraph{Withholding enforcement}
No protocol can attribute withholding to a specific actor, which suggests that slashing for omission faults is harder than slashing for commission faults such as equivocation or invalid execution.
Other enforcement models may be needed, for example insurance schemes, reputation systems with gradual stake reduction, or cryptographic mechanisms that make withholding unprofitable without having to attribute it.

\paragraph{Retention}
Four protocols define an expiry window, and Avail and Espresso leave retention to node operators.
This is by design: keeping all history indefinitely does not scale, and a DA layer exists to provide a window for verification and fraud detection.
Rollups and applications that need data beyond that window must store it elsewhere, for example on IPFS, Filecoin or Arweave.

\paragraph{Redundancy and commitments}
Among these six protocols, higher redundancy goes with simpler commitments.
Celestia (4$\times$) uses NMTs and Polkadot and NEAR (3$\times$) use Merkle trees, while Ethereum and Avail (2$\times$) use KZG, which verifies encoding without fraud proofs.
Espresso's polynomial VID also relies on algebraic commitments, which provide encoding verification and bribery resistance at a configurable redundancy.

\subsection{Limitations and Threats to Validity}
\label{sec:disc-limitations}

\paragraph{Demand-constrained measurements}
Observed throughput reflects demand during one 90-day window; conclusions about maximum throughput must come from protocol parameters (Section~\ref{sec:meth-normalization}).

\paragraph{Token prices}
USD costs depend on token prices, which moved widely during the window (ETH roughly \$2,000--\$4,000, TIA \$0.30--\$0.85), so another window would give different figures.
Reporting VWAP and percentiles reduces this effect but does not remove it.

\paragraph{Polkadot throughput upper bound}
Polkadot throughput is an upper bound because PoV sizes are not exposed on-chain; parachains typically use 1--3~MiB of the 10~MiB maximum.

\paragraph{Espresso Mainnet~0}
Mainnet~0 is a permissioned early deployment, so its cost, throughput and security figures cannot be compared on equal terms with permissionless, economically secured networks.
We include it for architectural completeness.

\paragraph{Single observation window}
A 90-day window is a snapshot, and behaviour under other market conditions, governance events or adoption levels may differ.
The released tooling can be used to extend the window.

\paragraph{Cross-protocol normalization}
Converting blob gas, coretime, gas and fixed fees to \$/MiB hides some economic detail: Polkadot's coretime also buys execution, and NEAR's gas price covers general contract execution.
The retention-adjusted cost (Table~\ref{tab:annualized_cost}, Fig.~\ref{fig:landscape}(b)) addresses part of this, but no single metric captures all of it.
